\documentclass[11pt,a4paper]{article}
\usepackage[margin=2.2cm]{geometry}
\usepackage{amsmath,amssymb,xcolor,enumitem,booktabs}
\usepackage[normalem]{ulem}
\usepackage[hidelinks]{hyperref}
\setlength{\parskip}{0.6em}\setlength{\parindent}{0pt}
\definecolor{newc}{HTML}{1f6f43}
\newcommand{\new}[1]{\textcolor{newc}{#1}}
\newcommand{\old}[1]{\textcolor{gray}{\sout{#1}}}
\title{Proposed change, Choosing together}
\date{20 November 2026, version 1}
\begin{document}\sloppy
\maketitle

\section*{How to read this}
This proposal applies the points of \texttt{method\_controllers\_comments.md}
to \texttt{sec:method:controllers} (pages 40 and 41 of \texttt{main.pdf}).
Part~1 is the subsection as it would read. Black text stays, \new{green} is
new and \old{struck} text goes. Part~2 lists what changes elsewhere, and
Part~3 the questions. Where a point needs your choice, the text below uses my
suggestion and the question says so. Nothing below is in the thesis yet.

Each change carries the point it answers, in brackets, for example [A1].

\section*{Part 1. The subsection}

\subsection*{D. Choosing together}

The centralised and section-based controllers use the same network, but on a
larger window and for one agent after another. The Problem Formulation gives
the order and the rule that keeps the agents of one decision apart
(Section V-F). \old{This subsection gives how the network carries them
out.} [D4]

\textbf{Centralised.} The network runs on \old{all vertices}\new{every
vertex an agent can stand on} [A3], once for each agent in priority order.
Each run uses the agent's own distance label and the vertices the earlier
agents have chosen, and the move it draws becomes part of what the next
agent sees. A timestep therefore takes $m$ runs, one after another, so
runtime grows with the size of the fleet. One pass shared by all agents, as
in \new{the single-agent setup of} Knippenberg, Holenderski, and Menkovski
(2021), would be cheaper, but then no agent would see another's choice. [A4]

\textbf{The claim mask.} The rule that keeps the agents of one decision
apart enters as a second mask\new{, the \emph{claim mask} $M^{\prec}_i$
(C.6). It switches off every move that rule forbids.}\old{. It switches off
a move to $w$ when an earlier agent has chosen $w$, when $w$ holds an agent
that has not chosen yet, or when the earlier agent standing on $w$ has
chosen to move to $v$, which would be a swap.} [D1] Waiting is never
switched off, so an agent always has an action left. \old{This is the claim
mask $M^{\prec}_i$ (C.6), and it}\new{The claim mask} is zero for a
decentralised agent. Conflict resolution therefore never overrides a
centralised agent\old{, and the override term of the reward (48) is always
zero for it}. [D1]

\textbf{\old{Seeing the whole window}\new{What the larger window adds}.}
[B1] After $L$ rounds a vertex only knows about vertices \new{at most} $L$
\old{hops}\new{edges} away, so a larger window alone adds little. The
centralised and section-based networks therefore add the mean vector over
the vertices the controller owns. \new{These are every vertex an agent can
stand on for the centralised controller, and the zone for the
section-based one.} [B2] This mean is also how an agent learns of earlier
choices more than $L$ \old{hops}\new{edges} away. \new{It is the readout
over all vertices of Gilmer et al.\ (2017) in its simplest form. A mean has
the same size for a window of any size, and the numbering of the vertices
does not change it.} [B2] It does not say which agent chose which vertex,
so the network cannot use every detail of the earlier choices, but the
claim mask still stops every collision.

In a large warehouse one earlier choice changes the mean very little.
\old{So the centralised controller mostly learns from choices close to the
agent, and the claim mask keeps the agents apart.}\new{So a centralised
agent sees the agents near it in detail and the rest of the warehouse only
as an average. What sets it apart from a decentralised agent is mainly that
it knows the earlier choices near it, that the claim mask keeps it clear of
them, and that it shares the credit of the whole fleet
(Section VI-G).} [B1]

These controllers use no messages (Section V-G), so \old{the message term
is replaced by}\new{the mean takes the place of the message term in} \new{the
readout (45)}. Formally, with $V_i^{\mathrm{pool}}$ the vertices the
controller owns, the \emph{pooled readout} is
\begin{equation*}
  \bar h_i(t) = \frac{1}{|V_i^{\mathrm{pool}}|}
    \sum_{v \in V_i^{\mathrm{pool}}} h_v^{(L)}, \qquad
  z_i = h_{\ell_i(t)}^{(L)} \;\Vert\; \bar h_i(t) \;\Vert\; \delta_i(t) .
  \tag{47}
\end{equation*}

\textbf{Section-based.} Each zone does the same on its own part of the
warehouse. \new{The zones are blocks of neighbouring aisles, as equal in
size as the aisles allow. A zone is not a section of SKUs, and the two are
drawn apart.} [C1] Zone $Y_q$ runs the network for each of its agents in
priority order, on its own vertices together with a halo $H_q$ (C.3). The
halo is the strip of vertices outside the zone within field-of-view range
of it. \old{Its vertices show their role and whether an agent stands there,
but they get no move scores and no loss, and they are left out of the mean
because they are only context.}\new{Its vertices carry the same facts as any
other vertex. No agent on them chooses in this zone, and they are left out
of the mean, because they are only context.} [A1] \old{The halo uses the
same range as the field of view, so the crowding feature means the same
under every controller.}\new{The halo reaches as far as the field of view,
so the window of an agent at the zone edge still holds its whole field of
view. A section-based agent therefore never sees a smaller part of the floor
than a decentralised one.} [A2] \new{The halo is our own construction.}
[B3]

A zone does not know the choices made in other zones, so an agent treats
every occupied halo vertex as holding an agent that has not chosen yet, and
does not move onto it. \old{An agent belongs to the next zone from the
timestep it enters it.} [D2] The weights are shared, so \old{this}\new{moving
into the next zone} changes the agent's context but not its policy. The
zones run in parallel, and two zones may still send agents onto the same
empty vertex at the border. Conflict resolution removes these clashes as it
does everywhere else. \new{So, unlike a centralised agent, a section-based
agent can be overridden at a zone border. This is part of what deciding in
zones costs, and the comparison keeps it.} [B4] \old{With $p=1$ the halo is
empty and the section-based controller is the centralised one.} [D2]

The three learned controllers share the network, the masks, the reward
\new{of each agent} and the training. They differ only in which agents
choose together, and so in the window the network runs on\old{ and}\new{,} in
which earlier choices an agent sees\new{, and in the credit each agent gets}. [D3] \old{A difference in results between them is therefore
a difference in how many agents one decision covers, which is the contrast
RQ1 asks about.}\new{This is the contrast RQ1 asks about.} [D3]

\section*{Part 2. What changes elsewhere}
\begin{itemize}[nosep]
  \item \texttt{B.Parameters.tex}, \texttt{tab:modelparams}, the row for
        the zones $p$. ``TBD (sweep)'' becomes 3 for the trained runs, with
        the scaling probe in its own appendix. [C1]
  \item \texttt{A.Notation.tex}. No change, $V_i^{\mathrm{pool}}$ and
        $\bar h_i(t)$ stay.
  \item \texttt{sec:method:rl}. The sentence that the override term is
        always zero for a centralised agent moves there, next to the term
        itself. [D1]
  \item \texttt{GAPS.md}. One new entry for the halo sentences and the zone
        cut, since both change what the text says the model does.
  \item The Discussion gets two points, what the centralised agent really
        sees [B1] and the overrides at zone borders [B4].
  \item No figure yet. [E1] waits for question~5.
\end{itemize}

\section*{Part 3. Questions}
\begin{enumerate}
  \item \textbf{Zone and section [C1].} I kept both names and said they
        differ. The other options are to rename the controller
        (zone-based) or the SKU sections. Which one?
  \item \textbf{What the centralised agent sees [B1].} Is the new
        paragraph fair to the centralised controller, or too strong?
  \item \textbf{The mean [B2].} I gave two reasons, the same size for any
        window and no dependence on numbering. Are those your reasons, or
        was there another one?
  \item \textbf{The halo [B3].} GAPS [A6] calls it the standard way to run a
        GNN on a split graph, but gives no source. I wrote ``our own''. Do
        you have a paper to cite instead?
  \item \textbf{A figure [E1].} Draw one zone, its halo and the three
        claim mask cases with $a_1$, $a_2$, in the viewer style?
  \item \textbf{Overrides at the border [B4].} Keep them as a cost of
        deciding in zones, as written?
\end{enumerate}

\end{document}
